\subsection{Galois cohomology}

Let N be a field, \(\Gamma\) a finite group of automorphisms of N and K the field of invariants of \(\Gamma\) . For every integer \(n \geqslant 1\) we denote by \(\mathbf{GL}(n, \mathbf{N})\) the group of square matrices of order n with coefficients in N and non-zero determinant (II, p.~349). We let the group \(\Gamma\) operate on the group \(\mathbf{GL}(n, \mathbf{N})\) by the rule \(\sigma(A) = (\sigma(a_{ij}))\) for \(A = (a_{ij})\) .

PROPOSITION 9. --- Let \((U_{\sigma})_{\sigma \in \Gamma}\) be a family of elements of \(\mathbf{GL}(n, \mathbf{N})\) . For \(A\) to exist in \(\mathbf{GL}(n, \mathbf{N})\) such that \(U_{\sigma} = A^{-1} \cdot \sigma(A)\) for all \(\sigma \in \Gamma\) it is necessary and sufficient that \(U_{\sigma\tau} = U_{\sigma} \cdot \sigma(U_{\tau})\) for \(\sigma, \tau\) in \(\Gamma\) .

The condition is necessary: if \(U_{\sigma} = A^{-1} \cdot \sigma(A)\) , then we have

\[
U _ {\sigma} \cdot \sigma (U _ {\tau}) = A ^ {- 1} \cdot \sigma (A) \sigma (A ^ {- 1} \cdot \tau (A)) = A ^ {- 1} \sigma \tau (A) = U _ {\sigma \tau}.
\]

The condition is sufficient: we identify the elements of \(\mathbf{N}^n\) with matrices of \(n\) rows and one column with coefficients in \(\mathbf{N}\) . We let the groups \(\Gamma\) act on \(\mathbf{N}^n\) by

\[
\sigma (x) = \left(\sigma \left(x _ {i}\right)\right) _ {1 \leqslant i \leqslant n} \quad \text { for } \quad x = \left(x _ {i}\right) _ {1 \leqslant i \leqslant n}.
\]

For each \(\sigma \in \Gamma\) we denote by \(u_{\sigma}\) the mapping \(x \mapsto U_{\sigma} \cdot \sigma(x)\) of \(\mathbf{N}^n\) into itself. The verification of Formulae (1) to (3) of V, p.~62 is immediate. Moreover we have \(u_{\varepsilon} \circ u_{\varepsilon} = u_{\varepsilon}\) and since \(u_{\varepsilon}\) is bijective, we have \(u_{\varepsilon} = \mathrm{Id}_{\mathbf{N}^n}\) . Let \(\mathbf{V}_0\) be the set of vectors \(x \in \mathbf{N}^n\) such that \(u_{\sigma}(x) = x\) for all \(\sigma \in \Gamma\) . By Prop. 7 (V, p.~63), \(\mathbf{V}_0\) is a K-structure on \(\mathbf{N}^n\) ; in particular there exist in \(\mathbf{V}_0\) vectors \(b_1, \ldots, b_n\) forming a basis of \(\mathbf{N}^n\) over N. The matrix B with columns \(b_1, \ldots, b_n\) is therefore invertible and the relation \(u_{\sigma}(b_i) = b_i\) for \(1 \leqslant i \leqslant n\) is equivalent to \(U_{\sigma} \cdot \sigma(B) = B\) . Writing \(A = B^{-1}\) , we obtain \(U_{\sigma} = A^{-1}\sigma(A)\) for all \(\sigma \in \Gamma\) .

COROLLARY 1. --- Let \((c_{\sigma})_{\sigma \in \Gamma}\) be a family of non-zero elements of N. For \(a \neq 0\) to exist in N such that \(c_{\sigma} = \sigma(a) \cdot a^{-1}\) for all \(\sigma \in \Gamma\) it is necessary and sufficient that \(c_{\sigma\tau} = c_{\sigma} \cdot \sigma(c_{\tau})\) for \(\sigma, \tau\) in \(\Gamma\) .

COROLLARY 2. --- Let \((c_{\sigma})_{\sigma \in \Gamma}\) be a family of elements of N. For b to exist in N such that \(a_{\sigma} = \sigma(b) - b\) for all \(\sigma \in \Gamma\) it is necessary and sufficient that \(a_{\sigma\tau} = a_{\sigma} + \sigma(a_{\tau})\) for \(\sigma, \tau\) in \(\Gamma\) .

We have \(\sigma \tau(b) - b = [\sigma(b) - b] + \sigma[\tau(b) - b]\) for all \(b\) in \(\mathbf{N}\) and \(\sigma, \tau\) in \(\Gamma\) , whence the necessity.

Conversely suppose that \(a_{\sigma \tau} = a_{\sigma} + \sigma(a_{\tau})\) for any \(\sigma\) and \(\tau\) in \(\Gamma\) . Put \(U_{\sigma} = \begin{pmatrix} 1 & a_{\sigma} \\ 0 & 1 \end{pmatrix}\) for \(\sigma \in \Gamma\) ; then we have \(U_{\sigma \tau} = U_{\sigma} \cdot \sigma(U_{\tau})\) for \(\sigma, \tau\) in \(\Gamma\) ; by Prop. 9, there exists thus a matrix \(A = \begin{pmatrix} x & y \\ z & t \end{pmatrix}\) with non-zero determinant such that \(\sigma(A) = AU_{\sigma}\) for all \(\sigma \in \Gamma\) ; writing down the relation \(\sigma(A) = AU_{\sigma}\) we find

\[
\left( \begin{array}{c c} \sigma (x) & \sigma (y) \\ \sigma (z) & \sigma (t) \end{array} \right) = \left( \begin{array}{c c} x & x a _ {\sigma} + y \\ z & z a _ {\sigma} + t \end{array} \right) (\sigma \in \Gamma).
\]

In particular, \(x\) and \(z\) belong to \(K\) and we have

\[
\sigma (y) = x a _ {\sigma} + y, \quad \sigma (t) = z a _ {\sigma} + t \quad (\sigma \in \Gamma).
\]

If \(x \neq 0\) we have \(a_{\sigma} = \sigma(b) - b\) with \(b = x^{-1}y\) ; if \(z \neq 0\) , we have the same relation with \(b = z^{-1}t\) . Now \(x\) and \(z\) cannot both be zero because

\[
x t - y z = \det A \neq 0.
\]

